\section{Research question and attribution}

The economic value of an exercise algorithm is the expected reward obtained by its stopping rule. Regression fit, numerical resolution and the number of simulated paths are means of constructing or evaluating that rule; none is itself a measure of value. A useful comparison must identify the financial problem, specify a feasible policy and state what its benchmark establishes. A finely resolved price can be a valuable numerical reference without being a known optimum. A valid upper bound can remain too loose to establish that policy loss is economically small.

The framework asks three related questions. How much expected value does a fixed exercise instruction forgo within a specified model? How much computation is needed to measure that loss with useful precision? How does the answer change when the financial inputs or opportunity set change? These questions require different comparisons. Changing the martingale used for evaluation changes bound tightness. Changing a policy changes its expected reward. Changing the exercise calendar can change the optimum even when the original policy is left untouched.

The methodological foundations are established. \citet{longstaff2001} construct exercise policies through backward continuation regression. \citet{haugh2004} develop martingale upper bounds and the use of approximate value processes. The lognormal hinge expectation used below follows the analytical pricing foundation of \citet{black1973}. The derivations here apply these foundations to an explicitly funded cash-event model and organise the resulting policy comparisons. The dual inequality, conditional Jensen argument and finite-family inference are not presented as new general theorems.

This document is the mathematical component of a wider research project. Its scope is deliberately independent of any particular observed contract, fitted parameter, sample size or empirical result. It defines the evidence that a numerical and financial investigation would need, rather than asserting that the evidence has been obtained.
